
\documentclass[11pt,a4paper]{article}

\usepackage[T1]{fontenc}
\usepackage[utf8]{inputenc}
\usepackage[ngerman]{babel}
\usepackage[a4paper,margin=25mm]{geometry}
\usepackage{amsmath,amssymb}
\usepackage{helvet}
\usepackage{sfmath}
\renewcommand{\familydefault}{\sfdefault}
\usepackage{microtype}
\usepackage[hidelinks]{hyperref}
\usepackage{lastpage}
\usepackage{fancyhdr}
\usepackage{parskip}

\pagestyle{fancy}
\fancyhf{}
\fancyfoot[C]{\thepage\ (\pageref{LastPage})}
\renewcommand{\headrulewidth}{0pt}
\renewcommand{\footrulewidth}{0pt}

\hypersetup{
  pdftitle={MoE-Glossar: DQP / Dichte QUBLE-Packung},
  pdfauthor={Wolfgang Tscheu}
}

\begin{document}

\vspace*{-1.8cm}

\noindent
\begin{minipage}[t]{0.66\textwidth}
{\LARGE\bfseries DQP / Dichte QUBLE-Packung}
\end{minipage}
\hfill
\begin{minipage}[t]{0.30\textwidth}
\raggedleft
{\small MoE-Glossareintrag}\\
{\small Version 261029-1}
\end{minipage}

\vspace{1.4em}

\section*{Definition}

Die Dichte QUBLE-Packung (DQP) ist eine
räumliche Ordnung gleich großer QUBLES.

Die QUBLES bilden durch gegenseitige
Oberflächenkontakte einen zusammenhängenden
Packungsverband.

Die räumliche Anordnung ihrer Kontaktlagen
bestimmt die Geometrie der Packung und
die Ausbildung der Interstitien.

\section*{Entstehung}

Die DQP geht aus der Umordnung der
Superdichten QUBLE-Packung (SDQP) hervor.

Im Verlauf der Expansion verändern sich
die Größenverhältnisse, Verformungen und
Kontaktlagen der QUBLES.

Beim Erreichen der geometrischen
Übergangsbedingung ordnet sich der
superdichte Packungsverband zur DQP um.

\section*{Interstitien}

Die räumliche Kontaktordnung der QUBLES
erzeugt tetraederähnliche und
oktaederähnliche Interstitien.

Vier einander paarweise berührende QUBLES
begrenzen ein konkaves, tetraederähnliches
Interstitium.

Sechs entsprechend angeordnete QUBLES
begrenzen ein konkaves, oktaederähnliches
Interstitium.

Dabei bilden vier QUBLE-Mittelpunkte
ein Quadrat. Je ein weiterer Mittelpunkt
liegt oberhalb und unterhalb der
Quadratsebene.

Die gekrümmten Oberflächen der QUBLES
bilden die Begrenzungsflächen der
Interstitien.

Die geometrischen Grundformen werden
im Glossareintrag \emph{Interstitium}
beschrieben.

\section*{Packungsordnung und Pfade}

Die DQP bildet eine räumliche
Kontaktordnung, innerhalb derer die
Pfadstrukturen des MoE beschrieben werden.

Die Geometrie der QUBLES und ihrer
Interstitien bestimmt die räumlichen
Bedingungen dieser Pfade.

\section*{Emergenzfolge}

\[
\text{Apeiron}
\rightarrow
\text{Apeiron-Expansion}
\rightarrow
\text{Quantelung}
\rightarrow
\text{QUBLE}
\rightarrow
\text{SDQP}
\rightarrow
\text{DQP}
\rightarrow
\text{Pfad}
\]

\vspace{0.9em}
\hrule
\vspace{0.45em}

{\footnotesize
MoE-Glossar:
\emph{QUBLE / Quantelung},
\emph{SDQP / Superdichte QUBLE-Packung},
\emph{Interstitium},
\emph{Pfad}.
}

\end{document}
